\subsection{Jordan Chains and Nilpotent Blocks}

Jordan chains convert nilpotent primary components into explicit blocks; each block records an eigenvalue and a chain length.

\begin{definition} \label{def:jordan-chain-block-and-matrix} \hlt{Jordan Chains, Blocks, and Matrices}
\begin{enumerate}[label=(\roman*)]
\setlength{\itemsep}{0pt}
\item Let $n,k\in\mathbb{Z}_{>0}$, $A\in\operatorname{Mat}(n,n,\mathbb{C})$, and $\lambda\in\mathbb{C}$. A \hlt{Jordan chain} for $A$ at $\lambda$ is $(\mathbf{x}_1,\ldots,\mathbf{x}_k)\in(\mathbb{C}^n)^k$ with $\mathbf{x}_1\neq\mathbf{0}$.\\
For $\mathbf{x}_0\defeq\mathbf{0}$, its defining relations are $(\forall j\in\{1,\ldots,k\})\,(A-\lambda I_n)\mathbf{x}_j=\mathbf{x}_{j-1}$.
\item For $k\in\mathbb{Z}_{>0}$, the \hlt{nilpotent Jordan block} $J_k(0)\in\operatorname{Mat}(k,k,\mathbb{C})$ satisfies $(\forall i,j\in\{1,\ldots,k\})\,[J_k(0)]_{ij}=\delta_{j,i+1}$.\\
For $\lambda\in\mathbb{C}$, the \hlt{Jordan block} at $\lambda$ is $J_k(\lambda)\defeq\lambda I_k+J_k(0)$.
\item A \hlt{Jordan matrix} is a finite direct sum of Jordan blocks.
\item Let $\mathcal{V}$ be a finite-dimensional complex vector space and let $T\in\LL(\mathcal{V})$.\\
An ordered basis $\beta$ of $\mathcal{V}$ is a \hlt{Jordan basis} for $T$ iff $[T]_{\beta\leftarrow\beta}$ is a Jordan matrix.
\end{enumerate}
\end{definition}

\begin{proposition} \label{prop:jordan-chains-as-matrix-columns} \hlt{Jordan Chains as Matrix Columns}\\
Let $n\in\mathbb{Z}_{>0}$ and let $A\in\operatorname{Mat}(n,n,\mathbb{C})$. Then:
\begin{enumerate}[label=(\roman*)]
\setlength{\itemsep}{0pt}
\item Let $k\in\mathbb{Z}_{>0}$ and $\lambda\in\mathbb{C}$, and let $(\mathbf{x}_1,\ldots,\mathbf{x}_k)$ be a Jordan chain for $A$ at $\lambda$. Set $X\defeq[\mathbf{x}_1\ \cdots\ \mathbf{x}_k]$; then $AX=XJ_k(\lambda)$.
\item Let $M\in\operatorname{Mat}(n,n,\mathbb{C})^{\times}$ be formed by concatenating Jordan chains for $A$, and let $J$ be the matching ordered direct sum of blocks.\\
Then $AM=MJ$ and $M^{-1}AM=J$.
\end{enumerate}
\begin{proof}
\begin{enumerate}[label=(\roman*)]
\setlength{\itemsep}{0pt}
\item Fix $j\in\{1,\ldots,k\}$. Then $(AX)_{\cdot,j}=\mathbf{x}_{j-1}+\lambda\mathbf{x}_j$ by the chain relation.\\
The identity $(XJ_k(\lambda))_{\cdot,j}=\mathbf{x}_{j-1}+\lambda\mathbf{x}_j$ holds by the entries of $J_k(\lambda)$. Hence $AX=XJ_k(\lambda)$ by equality of corresponding columns.
\item Write $M=[X_1\ \cdots\ X_s]$, where $X_i$ is the column matrix of a chain of length $k_i$ at $\lambda_i$.\\
The identities $(\forall i\in\{1,\ldots,s\})\,AX_i=X_iJ_{k_i}(\lambda_i)$ hold by part (i).\\
Thus $AM=MJ$ by block multiplication, and $M^{-1}AM=J$ by left multiplication with $M^{-1}$.
\end{enumerate}
\hlt{Structural Commentary}: Column identities pass from individual chains to their ordered block concatenation.\\
\hlt{Key Insight}: Superdiagonal ones encode predecessor vectors in successive columns.
\end{proof}
\end{proposition}

Figure \ref{fig:jordan-string-columns-similarity} depicts $N\mathbf{x}_1=\mathbf{0}$ and $N\mathbf{x}_j=\mathbf{x}_{j-1}$ along a Jordan chain.

\begin{illustrationfigure}[H]
  \centering
  \begin{tikzpicture}[
  >=Stealth,
  every node/.style={font=\small},
  vec/.style={draw=blue!65!black, fill=blue!7, rounded corners=2pt,
    minimum height=0.55cm, minimum width=0.95cm, align=center},
  arrow/.style={->, blue!70!black, line width=0.75pt},
  note/.style={text=gray!70!black, font=\footnotesize}
]
  \node[vec] (x1) at (0,0) {$\mathbf{x}_1$};
  \node[vec] (x2) at (1.63,0) {$\mathbf{x}_2$};
  \node[vec] (dots) at (3.21,0) {$\cdots$};
  \node[vec] (xk) at (4.80,0) {$\mathbf{x}_k$};

  \draw[arrow] (x1) -- ++(-0.90,0) node[left] {$\mathbf{0}$};
  \draw[arrow] (x2) -- node[above] {$N$} (x1);
  \draw[arrow] (dots) -- node[above] {$N$} (x2);
  \draw[arrow] (xk) -- node[above] {$N$} (dots);
\end{tikzpicture}

  \caption{A Jordan chain for $N=A-\lambda I$; each arrow records the image under $N$.}
  \label{fig:jordan-string-columns-similarity}
\end{illustrationfigure}

\begin{theorem} \label{thm:nilpotent-jordan-chain-basis} \hlt{Nilpotent Jordan Chain Basis}\\
Let $\mathcal{V}$ be a finite-dimensional vector space over $\mathbb{F}$,  $N\in\LL(\mathcal{V})$ be nilpotent. For $\mathbf{v}\in\mathcal{V}$ and $m\in\mathbb{Z}_{\geq0}$, set $\mathcal{C}_N(\mathbf{v},m)\defeq(N^{m-q}\mathbf{v})_{q=0}^{m}$.\\
For $r\in\mathbb{Z}_{\geq0}$ and $(\mathbf{v},\mathbf{m})\in\mathcal{V}^r\times\mathbb{Z}_{\geq0}^r$, let $\mathcal{B}_N(\mathbf{v},\mathbf{m})$ denote the concatenation $(\mathcal{C}_N(\mathbf{v}_j,m_j))_{j=1}^{r}$.\\
Write $\mathsf{CB}_N(r,\mathbf{v},\mathbf{m})$ when $\mathcal{B}_N(\mathbf{v},\mathbf{m})$ is a basis of $\mathcal{V}$ and $(\forall j\in\{1,\ldots,r\})\,N^{m_j+1}\mathbf{v}_j=\mathbf{0}$.\\
Then $(\exists r\in\mathbb{Z}_{\geq0})(\exists\mathbf{v}\in\mathcal{V}^r)(\exists\mathbf{m}\in\mathbb{Z}_{\geq0}^r)\,\mathsf{CB}_N(r,\mathbf{v},\mathbf{m})$.
\begin{proof}
Set $n\defeq\dim(\mathcal{V})$ and argue by strong induction on $n$, assuming the claim in every smaller dimension.
\begin{enumerate}[label=(\roman*)]
\setlength{\itemsep}{0pt}
\item \hlt{Zero Operator}: If $N=0$, choose an ordered basis $(\mathbf{v}_1,\ldots,\mathbf{v}_r)$ by Corollary \ref{cor:finite-dimensional-spaces-and-subspaces-have-bases}(i).\\
Set $\mathbf{m}\defeq(0,\ldots,0)\in\mathbb{Z}_{\geq0}^r$; then $\mathcal{B}_N(\mathbf{v},\mathbf{m})=(\mathbf{v}_1,\ldots,\mathbf{v}_r)$. The terminal relations $(\forall j\in\{1,\ldots,r\})\,N\mathbf{v}_j=\mathbf{0}$ hold by $N=0$.\\
Thus $\mathsf{CB}_N(r,\mathbf{v},\mathbf{m})$ by its defining conditions; when $\mathcal{V}=\{\mathbf{0}\}$, $r=0$ by linear independence of the chosen basis.
\item \hlt{Range Reduction}: Assume $N\neq0$. Then $n>0$ and $N^n=0$ by Corollary \ref{cor:nilpotent-power-dimension-zero}.\\
If $N$ were injective, then $N^n$ would be injective by induction on $n$, contrary to $N^n=0$ on the nonzero space $\mathcal{V}$.\\
Thus $\ker(N)\neq\{\mathbf{0}\}$ by noninjectivity, and $\dim\operatorname{range}(N)<n$ by Theorem \ref{thm:rank-nullity}(ii).\\
Set $\mathcal{W}\defeq\operatorname{range}(N)$. The restriction $N|_{\mathcal{W}}$ is nilpotent by $N(\mathcal{W})\subseteq\mathcal{W}$ and $(N|_{\mathcal{W}})^n=0$.\\
Choose $t\in\mathbb{Z}_{\geq0}$, $\mathbf{u}=(\mathbf{u}_j)_{j=1}^{t}\in\mathcal{W}^t$, and $\boldsymbol{\ell}=(\ell_j)_{j=1}^{t}\in\mathbb{Z}_{\geq0}^t$ by the induction hypothesis for $N|_{\mathcal{W}}$.\\
Set $\mathcal{B}_{\mathcal{W}}\defeq(\mathcal{C}_N(\mathbf{u}_j,\ell_j))_{j=1}^{t}$. This is a basis of $\mathcal{W}$ by equality of restriction powers and the induction hypothesis.\\
The identities $(\forall j\in\{1,\ldots,t\})\,N^{\ell_j+1}\mathbf{u}_j=\mathbf{0}$ hold by the same facts.\\
For $j\in\{1,\ldots,t\}$, choose $\mathbf{v}_j\in\mathcal{V}$ satisfying $N\mathbf{v}_j=\mathbf{u}_j$ by $\mathbf{u}_j\in\operatorname{range}(N)$. Set $\mathcal{L}\defeq(\mathcal{C}_N(\mathbf{v}_j,\ell_j+1))_{j=1}^{t}$.
\item \hlt{Lifted Independence}: Let $c_{j,q}\in\mathbb{F}$, where $1\leq j\leq t$ and $0\leq q\leq\ell_j+1$.\\
Suppose $\sum\limits_{j=1}^{t}\sum\limits_{q=0}^{\ell_j+1}c_{j,q}N^q\mathbf{v}_j=\mathbf{0}$. The identities $(\forall j\in\{1,\ldots,t\})\,N^{\ell_j+2}\mathbf{v}_j=N^{\ell_j+1}\mathbf{u}_j=\mathbf{0}$ hold by item (ii).\\
The image of the relation under $N$ is $\sum\limits_{j=1}^{t}\sum\limits_{q=0}^{\ell_j}c_{j,q}N^q\mathbf{u}_j=\mathbf{0}$ by linearity and the terminal identities.\\
Thus $(\forall j\in\{1,\ldots,t\})(\forall q\in\{0,\ldots,\ell_j\})\,c_{j,q}=0$ by linear independence of $\mathcal{B}_{\mathcal{W}}$.\\
The original relation reduces to $\sum\limits_{j=1}^{t}c_{j,\ell_j+1}N^{\ell_j}\mathbf{u}_j=\mathbf{0}$ by these coefficient equalities.\\
Every remaining coefficient is zero by linear independence of $\mathcal{B}_{\mathcal{W}}$. Thus $\mathcal{L}$ is linearly independent by the vanishing of all coefficients.
\item \hlt{Kernel Completion}: The identities $(\forall j\in\{1,\ldots,t\})(\forall q\in\{0,\ldots,\ell_j\})\,N(N^q\mathbf{v}_j)=N^q\mathbf{u}_j$ hold by $N\mathbf{v}_j=\mathbf{u}_j$.\\
The inclusion $\mathcal{W}\subseteq N(\operatorname{span}(\mathcal{L}))$ holds by spanning of $\mathcal{B}_{\mathcal{W}}$. The reverse inclusion holds by $\mathcal{W}=\operatorname{range}(N)$, so $\mathcal{W}=N(\operatorname{span}(\mathcal{L}))$.\\
Choose $\mathbf{w}_1,\ldots,\mathbf{w}_s$ such that $(\mathcal{L},\mathbf{w}_1,\ldots,\mathbf{w}_s)$ is a basis by Theorem \ref{thm:extending-independent-lists}.\\
Choose $(\mathbf{x}_i)_{i=1}^{s}\in\operatorname{span}(\mathcal{L})^s$ satisfying $(\forall i\in\{1,\ldots,s\})\,N\mathbf{x}_i=N\mathbf{w}_i$ by $\mathcal{W}=N(\operatorname{span}(\mathcal{L}))$.\\
Set $\mathbf{z}_i\defeq\mathbf{w}_i-\mathbf{x}_i$ for $i\in\{1,\ldots,s\}$. Then $(\forall i\in\{1,\ldots,s\})\,N\mathbf{z}_i=\mathbf{0}$ by construction.\\
The list $(\mathcal{L},\mathbf{z}_1,\ldots,\mathbf{z}_s)$ spans $\mathcal{V}$ by $\mathbf{w}_i=\mathbf{z}_i+\mathbf{x}_i$ and $\mathbf{x}_i\in\operatorname{span}(\mathcal{L})$. Suppose a linear combination of $(\mathcal{L},\mathbf{z}_1,\ldots,\mathbf{z}_s)$ is zero.\\
Substitute $\mathbf{z}_i=\mathbf{w}_i-\mathbf{x}_i$ in this relation. Every $\mathbf{w}_i$-coefficient is zero by linear independence of $(\mathcal{L},\mathbf{w}_1,\ldots,\mathbf{w}_s)$.\\
The remaining $\mathcal{L}$-coefficients are zero by linear independence of $\mathcal{L}$.\\
Thus $(\mathcal{L},\mathbf{z}_1,\ldots,\mathbf{z}_s)$ is linearly independent, hence a basis by the preceding spanning property.
\item \hlt{Chain Verification}: Set $r\defeq t+s$ and $\mathbf{v}'\defeq(\mathbf{v}_1,\ldots,\mathbf{v}_t,\mathbf{z}_1,\ldots,\mathbf{z}_s)$. Set $\mathbf{m}'\defeq(\ell_1+1,\ldots,\ell_t+1,\underbrace{0,\ldots,0}_{s\text{ entries}})$.\\
Then $\mathcal{B}_N(\mathbf{v}',\mathbf{m}')=(\mathcal{L},\mathbf{z}_1,\ldots,\mathbf{z}_s)$ is a basis by item (iv). The identities $(\forall j\in\{1,\ldots,t\})\,N^{\ell_j+2}\mathbf{v}_j=\mathbf{0}$ hold by item (iii).\\
The identities $(\forall i\in\{1,\ldots,s\})\,N\mathbf{z}_i=\mathbf{0}$ hold by item (iv). Thus $\mathsf{CB}_N(r,\mathbf{v}',\mathbf{m}')$ by its defining conditions.
\end{enumerate}
\hlt{Structural Commentary}: Strong induction lifts a range-chain basis; kernel correction completes the lifted list.\\
\hlt{Key Insight}: The $N$-image of a lifted dependence lies in the known range basis.
\end{proof}
\end{theorem}

Figure \ref{fig:jordan-chain-nilpotent-basis} illustrates the nilpotent shift along one chain in the adapted basis order.

\begin{illustrationfigure}[H]
  \centering
  \begin{tikzpicture}[
  font=\small,
  vec/.style={draw=blue!65!black, fill=blue!7, rounded corners=2pt, minimum height=0.58cm, minimum width=1.12cm, align=center},
  arrow/.style={-{Stealth[length=2mm]}, thick, blue!70!black}
]
  \node[vec] (a0) at (0,0) {$N^m\mathbf{v}$};
  \node[vec] (a1) at (1.76,0) {$N^{m-1}\mathbf{v}$};
  \node[vec] (a2) at (3.34,0) {$\cdots$};
  \node[vec] (a3) at (4.88,0) {$N\mathbf{v}$};
  \node[vec] (a4) at (6.38,0) {$\mathbf{v}$};

  \draw[arrow] (a4) -- node[above] {$N$} (a3);
  \draw[arrow] (a3) -- node[above] {$N$} (a2);
  \draw[arrow] (a2) -- node[above] {$N$} (a1);
  \draw[arrow] (a1) -- node[above] {$N$} (a0);
  \draw[arrow] (a0) -- ++(-1.18,0) node[left] {$\mathbf{0}$};
\end{tikzpicture}

  \caption{An ordered nilpotent chain shifts left under $N$ and represents one Jordan block.}
  \label{fig:jordan-chain-nilpotent-basis}
\end{illustrationfigure}

\begin{computationmethod} \label{meth:recursive-range-lifting-jordan-construction} \hlt{Recursive Range-Lifting Jordan Construction}\\
Given $n\in\mathbb{Z}_{>0}$ and $A\in\operatorname{Mat}(n,n,\mathbb{C})$, construct basis matrix $M\in\operatorname{Mat}(n,n,\mathbb{C})^{\times}$. Construct Jordan matrix $J$ with postcondition $AM=MJ$.\\
Fix $\lambda\in\operatorname{spec}(A)$ and set $\mathcal{G}_{\lambda}\defeq\ker((A-\lambda I_n)^n)$. Set $N_{\lambda}\defeq(A-\lambda I_n)|_{\mathcal{G}_{\lambda}}$.\\
The inclusion $(A-\lambda I_n)\mathcal{G}_{\lambda}\subseteq\mathcal{G}_{\lambda}$ holds by commutativity of powers of $A-\lambda I_n$. Thus $N_{\lambda}\in\LL(\mathcal{G}_{\lambda})$ and $N_{\lambda}^n=0$ by the definition of $\mathcal{G}_{\lambda}$.
\begin{enumerate}[label=(\roman*)]
\setlength{\itemsep}{0pt}
\item If $N_{\lambda}=0$, choose an ordered basis of $\mathcal{G}_{\lambda}$ and regard its vectors as singleton chains.
\item If $N_{\lambda}\neq0$, recursively construct a chain basis for $N_{\lambda}|_{\operatorname{range}(N_{\lambda})}$.
\item Given a returned chain $\mathcal{C}_{N_{\lambda}}(\mathbf{u},\ell)$, choose $\mathbf{v}\in\mathcal{G}_{\lambda}$ with $N_{\lambda}\mathbf{v}=\mathbf{u}$. Replace the returned chain by $\mathcal{C}_{N_{\lambda}}(\mathbf{v},\ell+1)$.
\item Perform the kernel-completion correction from the proof of Theorem \ref{thm:nilpotent-jordan-chain-basis} on the lifted list.\\
Append the resulting kernel vectors as singleton chains.
\item Repeat (i)--(iv) over $\lambda\in\operatorname{spec}(A)$, concatenate the resulting bases, and order the matching Jordan blocks identically.
\end{enumerate}
If $N_{\lambda}\neq0$, then $N_{\lambda}$ is not injective because injectivity would make $N_{\lambda}^n=0$ injective.\\
Thus its range dimension is strictly smaller by Theorem \ref{thm:rank-nullity}(ii), so the recursion terminates.\\
Steps (i)--(iv) produce a basis of $\mathcal{G}_{\lambda}$ by Theorem \ref{thm:nilpotent-jordan-chain-basis}.\\
The concatenation in (v) is a basis by Theorem \ref{thm:complex-primary-decompositi generalized-eigenspaces}(i).\\
Each returned $\mathcal{C}_{N_{\lambda}}(\mathbf{v},m)$ satisfies $N_{\lambda}^{m+1}\mathbf{v}=\mathbf{0}$ by construction. It is a Jordan chain for $A$ at $\lambda$ by $N_{\lambda}=A-\lambda I_n$ on $\mathcal{G}_{\lambda}$. The postcondition $AM=MJ$ holds by Proposition \ref{prop:jordan-chains-as-matrix-columns}(ii).
\end{computationmethod}

\begin{proposition} \label{prop:powers-of-jordan-block} \hlt{Powers of a Jordan Block}\\
Let $k\in\mathbb{Z}_{>0}$ and $\lambda\in\mathbb{C}$, using $0^0\defeq1$ in (i). Then:
\begin{enumerate}[label=(\roman*)]
\setlength{\itemsep}{0pt}
\item $(\forall\,p\in\mathbb{Z}_{\geq0})\,J_k(\lambda)^p=\sum\limits_{r=0}^{\min(p,k-1)}\binom{p}{r}\lambda^{p-r}J_k(0)^r$.
\item $J_k(\lambda)^p\to0$ entrywise as $p\to\infty$ iff $\abs{\lambda}<1$.
\end{enumerate}
\begin{proof}
\begin{enumerate}[label=(\roman*)]
\setlength{\itemsep}{0pt}
\item[$(i)$] Fix $p\in\mathbb{Z}_{\geq0}$, $R\defeq J_k(0)$. For $0\leq r<k$, matrix $R^r$ has ones on $r$th superdiagonal, zeros elsewhere by induction on $r$.\\
The identity $(\forall r\in\mathbb{Z}_{\geq k})\,R^r=0$ holds by the same induction.\\
By distributivity, each subset of the $p$ factors in $(\lambda I_k+R)^p$ selects the factors contributing $R$.\\
For $r\in\{0,\ldots,p\}$, exactly $\binom{p}{r}$ subsets select $r$ factors by Proposition \ref{prop:unordered-samples-binomial-coefficients}.\\
The identity $(\lambda I_k+R)^p=\sum\limits_{r=0}^{p}\binom{p}{r}\lambda^{p-r}R^r$ holds by commutativity and regrouping by subset size.\\
The identity in (i) holds by deleting the zero terms with $r\geq k$.
\item[$(ii)$] ($\Rightarrow$) The diagonal sequence satisfies $\lambda^p\to0$ by entrywise convergence of $J_k(\lambda)^p$.\\
If $\abs{\lambda}\geq1$, then $(\forall p\in\mathbb{Z}_{\geq0})\,\abs{\lambda^p}\geq1$, contrary to $\lambda^p\to0$. Thus $\abs{\lambda}<1$ by contradiction.
\item[$(ii)$] ($\Leftarrow$) If $\lambda=0$, then $(\forall p\in\mathbb{Z}_{\geq k})\,J_k(0)^p=0$ by the power description in (i).\\
Assume $\rho\defeq\abs{\lambda}\in(0,1)$ and fix $r\in\{0,\ldots,k-1\}$. Set $a_{p,r}\defeq\binom{p}{r}\rho^{p-r}$ for $p\geq r$.\\
Set $\theta\defeq\frac{1+\rho}{2}\in(\rho,1)$. Choose $P\in\mathbb{Z}_{\geq r}$ satisfying $(\theta-\rho)(P+1)\geq\theta r$ by unboundedness of $\mathbb{Z}_{>0}$.\\
$(\forall p\in\mathbb{Z}_{\geq P})\,\frac{a_{p+1,r}}{a_{p,r}}=\frac{\rho(p+1)}{p+1-r}\leq\theta$ by factorial cancellation and the choice of $P$.\\
Hence $(\forall m\in\mathbb{Z}_{\geq0})\,a_{P+m,r}\leq a_{P,r}\theta^m$ by induction on $m$. Fix $\varepsilon>0$ and set $\delta\defeq\theta^{-1}-1>0$.\\
Choose $M\in\mathbb{Z}_{>0}$ satisfying $1+M\delta>\frac{a_{P,r}}{\varepsilon}$ by unbounded $\mathbb{Z}_{>0}$. For $m\in\mathbb{Z}_{\geq M}$, $\theta^{-m}=(1+\delta)^m\geq1+m\delta$ by Proposition \ref{prop:binomial-theorem-and-subset-count}(i).\\
Thus $0\leq a_{P+m,r}\leq\frac{a_{P,r}}{1+m\delta}<\varepsilon$, so $a_{p,r}\to0$. Thus $(\forall r\in\{0,\ldots,k-1\})\,\binom{p}{r}\lambda^{p-r}\to0$ as $p\to\infty$.\\
Each entry in the expansion from (i) is zero or $\binom{p}{r}\lambda^{p-r}$ for one $r\in\{0,\ldots,k-1\}$.\\
Its absolute value is therefore either $0$ or $a_{p,r}$, and hence tends to $0$. Thus $J_k(\lambda)^p\to0$ by convergence of each entry.
\end{enumerate}
\hlt{Structural Commentary}: Nilpotence truncates the expansion; diagonal entries and a geometric majorant resolve the limit criterion.\\
\hlt{Key Insight}: Each fixed superdiagonal has polynomial coefficients multiplied by geometric decay.
\end{proof}
\end{proposition}
